\subsection{复指数函数}

我们已在（Gen.~Top., VIII, p.~106）中确定了复数的（加法）拓扑群 C 到（乘法）拓扑群 \(C^{*}\) （由 \(\neq 0\) 的复数组成）上的所有连续同态；它们是映射

\[
x + i y \mapsto e ^ {\alpha x + \beta y} \mathbf {e} (\gamma x + \delta y)\tag{20}
\]

其中 \(\alpha, \beta, \gamma, \delta\) 是满足唯一条件 \(\alpha\delta - \beta\gamma \neq 0\) 的四个实数。我们现在要确定这些同态 \(z \mapsto f(z)\) 中哪些在 C 上可导。首先注意，只需 f 在点 z = 0 可导即可；事实上，对每一点 \(z \in C\) 有 \(\frac{f(z + h) - f(z)}{h} = f(z) \frac{f(h) - 1}{h}\) ；若 \(f'(0)\) 存在，则 \(f'(z)\) 也存在，且 \(f'(z) = af(z)\) ，其中 \(a = f'(0)\) 。另一方面，若 g 是第二个可导同态，满足 \(g'(z) = bg(z)\) ，则 \(g(az/b) = f(z)\) ，因为立即可以看出商 \(g(az/b)/f(z)\) 的导数处处为零且在 z = 0 处等于 1；因此所有可导同态都具有 \(z \mapsto f(\lambda z)\) 的形式，其中 f 是它们中的某一个（假定它们存在），而 \(\lambda\) 是任意（复）常数。

在此情形下，若 \(f\) 在点 \(z = 0\) 可导，则映射 \(x \mapsto f(x)\) 与 \(y \mapsto f(iy)\) （都是 \(\mathbf{R}\) 到 \(\mathbf{C}\) 中的映射）都必在点 0 可导，前者的导数为 \(f'(0)\) ，后者的导数为 \(if'(0)\) 。而映射 \(x \mapsto e^{\alpha x}\mathbf{e}(\gamma x)\) 与 \(y \mapsto e^{\beta y}\mathbf{e}(\delta y)\) 在点 0 的导数分别等于 \(\alpha + 2\pi i\gamma\) 与 \(\beta + 2\pi i\delta\) ，由此得 \(\beta = -2\pi\gamma\) 与 \(\alpha = 2\pi\delta\) ；这些条件特别地为同态 \(x + iy \mapsto e^x\mathbf{e}(y/2\pi)\) 所满足，我们暂时把它记作 \(f_0\) 。现在证明 \(f_0\) 确实在点 \(z = 0\) 可导。

显然 \(x \mapsto f_0(x)\) 与 \(y \mapsto f_0(iy)\) 有各阶导数；特别地，对这两个函数应用 1 阶泰勒公式可知，对每个 \(\varepsilon > 0\) 存在 \(r > 0\) ，使得若令

\[
f _ {0} (x) = 1 + x + \varphi (x) x, \quad f _ {0} (i y) = 1 + i y + \psi (y) y,
\]

则条件 \(|x| \leqslant r, |y| \leqslant r\) 蕴含 \(|\varphi(x)| \leqslant \varepsilon\) 与 \(|\psi(y)| \leqslant \varepsilon\) ；在此情形下，有 \(f_0(x + iy) = f_0(x)f_0(iy) = 1 + (x + iy) + \theta(x, y)\) ，其中

\[
\theta (x, y) = \bigl (i + \varphi (x) \psi (y) \bigr) x y + (1 + x) y \psi (y) + (1 + i y) x \varphi (x);
\]

当 \(|z| \leqslant r\) 时有 \(|x| \leqslant r\) 与 \(|y| \leqslant r\) ，从而

\[
\left| \theta (x, y) \right| \leqslant \left(1 + \varepsilon^ {2}\right) | z | ^ {2} + 2 \varepsilon | z | (1 + | z |)
\]

这证明商 \(\frac{f_{0}(z)-1-z}{z}\) 随 z 趋于 0 而趋于 0 ，也就是说，函数 \(f_{0}\) 在点 z=0 处有导数，等于 1 。于是上面证明了，对一切 \(z\in C\) ，

\[
\mathrm{D} \big (f _ {0} (z) \big) = f _ {0} (z).\tag{21}
\]

这一性质建立了 \(f_{0}\) 与函数 \(e^{x}\) 之间的联系，后者正是 \(f_{0}\) 在实轴上的限制；为此我们给出下面的定义：

定义 2. 同态 \(x + iy \mapsto e^x\mathbf{e}(y/2\pi)\) （\(\mathbf{C}\) 到 \(\mathbf{C}^*\) 上的）称为复指数（函数）；它在任意复数 \(z\) 处的值记作 \(e^z\) 或 \(\exp z\) 。

